\maketitle
\begin{abstract}
\Lang{A difference between models can remain unknown, fail to be observable with present resources, or fail to produce any response under an entire declared future protocol class. These are different statements. We attach explicit time and energy budgets to observer protocols and give exact and bounded-error separation certificates, while retaining the response equivalence inherited from WR-I. A finite adaptive transcript cannot distinguish models whose conditional response laws agree at every common feasible history. The concrete benchmark is an ideal tagged light pulse in a flat expanding cubic torus, compared with its Euclidean cover. Its first return exists exactly when the torus period is strictly smaller than the remaining comoving light-travel span. A detector threshold adds a separate energy condition. For exponential expansion we derive the exact time--energy detection frontier, including its unattained causal endpoint. Smooth positive expansion histories can agree on the entire supplied past and on an arbitrarily long finite future, yet have finite versus infinite remaining conformal span. Thus current background compatibility does not determine permanent future accessibility without additional dynamical assumptions. The ninth paper closes the first technical cycle at the level of a reviewable mathematical draft; it establishes no universal limit on knowledge or new cosmological law.}{Различие моделей может оставаться неизвестным, быть недоступным при нынешних ресурсах или не порождать различимого отклика во всём объявленном классе будущих протоколов. Это разные утверждения. Мы вводим явные бюджеты времени и энергии для протоколов наблюдателя и даём сертификаты точного различения и различения при ограниченной ошибке, сохраняя эквивалентность откликов из WR-I. Конечная адаптивная запись не различает модели, условные законы откликов которых совпадают при каждой общей допустимой истории. Конкретный пример \textemdash{} идеальный меченый световой импульс в плоском расширяющемся кубическом торе в сравнении с его евклидовым накрытием. Первое возвращение существует тогда и только тогда, когда период тора строго меньше оставшегося сопутствующего пути света. Порог детектора добавляет отдельное энергетическое условие. Для экспоненциального расширения выведена точная граница обнаружения по времени и энергии с недостижимой причинной граничной точкой. Гладкие положительные истории расширения могут совпадать на всём заданном прошлом и сколь угодно длинном конечном будущем, но иметь конечный либо бесконечный оставшийся конформный путь. Поэтому совместимость с нынешним фоном не определяет постоянную будущую доступность без дополнительных динамических предпосылок. Девятая статья завершает первый технический цикл на уровне математического черновика для рецензирования; универсальный предел познания или новый космологический закон не устанавливается.}
\end{abstract}
\noindent\textbf{\Lang{Status}{Статус}:} REVIEWABLE\_DRAFT. \Lang{Independent external scientific peer review is not established. No WR-IX DOI has been assigned.}{Независимое внешнее научное рецензирование не установлено. DOI WR-IX не присвоен.}

\section{\Lang{The final obligation of the first cycle}{Последняя задача первого цикла}}\label{sec:contract}
\Lang{The WREH manifesto assigns WR-IX the distinction between unknown, currently unobservable, inaccessible in principle and physically indistinguishable \cite{Manifesto}. WR-I already owns response equivalence and finite separation on compact domains \cite{WRI}; WR-II--VI own existence, fibre geometry, historical and experimental completion, and coupled admissibility \cite{WRII,WRIII,WRIV,WRV,WRVI}. WR-VII supplies a calibrated unit-bearing quantum resource benchmark, rather than a universal observer resource bound \cite{WRVII}. WR-VIII supplies a finite current cosmological patch and alternative kinematic lifts \cite{WRVIII}.}{Манифест WREH закрепляет за WR-IX различие между неизвестным, сейчас ненаблюдаемым, принципиально недоступным и физически неразличимым \cite{Manifesto}. WR-I уже содержит эквивалентность откликов и конечное разделение на компактных областях \cite{WRI}; WR-II--VI посвящены существованию, геометрии слоёв, историческому и экспериментальному достраиванию и связанным условиям допустимости \cite{WRII,WRIII,WRIV,WRV,WRVI}. WR-VII даёт квантовый ресурсный пример с калиброванными физическими единицами, а не универсальную границу ресурсов наблюдателя \cite{WRVII}. WR-VIII задаёт конечный нынешний космологический участок и альтернативные кинематические продолжения \cite{WRVIII}.}

\Lang{The new obligation is to construct an observer-accessible future response family and its resource frontier, then test whether the frontier is determined by current data. We do not relabel the WR-I quotient as a new result, identify refinement order with physical time, or import the qubit carrier into cosmology. The mathematical arguments below are elementary and self-contained; no theorem-priority claim or exhaustive novelty certificate is made.}{Новая задача \textemdash{} построить доступное наблюдателю семейство будущих откликов и его ресурсную границу, а затем проверить, определяется ли эта граница нынешними данными. Мы не переименовываем фактор WR-I в новый результат, не отождествляем порядок уточнения с физическим временем и не переносим кубитный носитель в космологию. Математические аргументы ниже элементарны и самодостаточны; приоритет теорем и исчерпывающая новизна не заявляются.}

\section{\Lang{Access records and the scope of a horizon}{Условия доступа и область действия горизонта}}\label{sec:access}
\Lang{Fix a nonempty current compatibility class $C$, an observer/laboratory description $O$, and a class $\mathcal P$ of complete finite protocols. Each protocol has a common response map $r_p:C\to\mathbb R$ and certified upper bounds $t_p\ge0$, $e_p\ge0$ on elapsed proper time and launch energy. Units are seconds and joules. Scalar responses are normalized in stated response units. For this first abstract interface, feasibility and the certified bounds hold for every $W\in C$.}{Зафиксируем непустой нынешний класс совместимости $C$, описание наблюдателя и лаборатории $O$ и класс $\mathcal P$ полных конечных протоколов. Каждый протокол имеет общую карту отклика $r_p:C\to\mathbb R$ и сертифицированные верхние оценки $t_p\ge0$, $e_p\ge0$ прошедшего собственного времени и энергии запуска. Единицы \textemdash{} секунды и джоули. Скалярные отклики нормированы в указанных единицах отклика. В этом первом абстрактном интерфейсе выполнимость и оценки ресурсов действуют для каждого $W\in C$.}
\begin{equation}\label{eq:access}
 \mathcal A_{\tau,E}=\{p\in\mathcal P:t_p\le\tau,\ e_p\le E\},\qquad
 \mathcal A_\infty=\bigcup_{\tau,E<\infty}\mathcal A_{\tau,E}.
\end{equation}
\Lang{Here $\tau,E\ge0$; the union permits arbitrarily large finite resources, not an experiment completed at infinite time. Costs of a complete adaptive protocol must include every branch used in its certificate. Known gauge redundancies should already be identified in the declared carrier. The observer, model class, allowed preparations and cost convention are part of the access record, not inferred from a response quotient. If feasibility depends on $W$, taking the intersection of feasible labels over $C$ defines a conservative common class. Its horizon need not survive model-dependent feasible actions, learning, or an enlarged physical carrier. An informative failure must be represented as an outcome, not silently deleted.}{Здесь $\tau,E\ge0$; объединение допускает сколь угодно большие конечные ресурсы, но не эксперимент, завершённый в бесконечный момент. Стоимость полного адаптивного протокола должна учитывать все ветви его сертификата. Известные калибровочные избыточности следует заранее отождествить в объявленном носителе. Наблюдатель, класс моделей, допустимые подготовки и соглашение о стоимости входят в условия доступа, а не выводятся из фактора откликов. Если выполнимость зависит от $W$, пересечение допустимых меток по $C$ задаёт консервативный общий класс. Его горизонт может исчезнуть при действиях с зависящей от модели выполнимостью, обучении или расширении физического носителя. Информативный отказ следует включать в исходы, а не молча исключать.}
\begin{equation}\label{eq:gap}
 g_{\tau,E}(W,V)=\sup_{p\in\mathcal A_{\tau,E}}|r_p(W)-r_p(V)|,\qquad
 g_\infty(W,V)=\sup_{p\in\mathcal A_\infty}|r_p(W)-r_p(V)|.
\end{equation}
\Lang{Set the supremum of an empty protocol family to zero; infinite values are allowed. A pair with $W\ne V$ and $g_\infty=0$ has an exact horizon relative to this access record. This asserts response equality, not ontological identity. ``Unknown'' describes missing information; ``currently inaccessible'' refers to a specified finite budget; ``permanently inaccessible'' quantifies over the declared union. A missing certificate is not a proof of a horizon.}{Супремум пустого семейства протоколов полагаем равным нулю; бесконечные значения допускаются. Пара $W\ne V$ с $g_\infty=0$ имеет точный горизонт относительно этих условий доступа. Это равенство откликов, а не онтологическое тождество. «Неизвестно» описывает нехватку информации; «сейчас недоступно» относится к заданному конечному бюджету; «постоянно недоступно» квантифицирует по объявленному объединению. Отсутствие сертификата не доказывает горизонт.}

\section{\Lang{Separation, noise and adaptive protocols}{Различение, шум и адаптивные протоколы}}\label{sec:separation}
\begin{proposition}[\Lang{Finite witnesses and a strict error threshold}{Конечные свидетельства и строгий порог ошибки}]\label{prop:witness}
\Lang{For the access record above,}{Для заданных выше условий доступа}
\begin{equation}\label{eq:sup}
 g_\infty=\sup_{\tau,E<\infty}g_{\tau,E}.
\end{equation}
\Lang{A pair is exactly separated by some finite protocol if and only if $g_\infty>0$. Suppose the complete scalar observation has a deterministic additive error in $[-\varepsilon,\varepsilon]$, with the same $\varepsilon\ge0$ for each protocol. A single-protocol worst-case certificate exists if and only if $g_\infty>2\varepsilon$.}{Пара точно различается некоторым конечным протоколом тогда и только тогда, когда $g_\infty>0$. Пусть полное скалярное наблюдение имеет детерминированную аддитивную ошибку в $[-\varepsilon,\varepsilon]$ с одним и тем же $\varepsilon\ge0$ для каждого протокола. Сертификат одного протокола для худшего случая существует тогда и только тогда, когда $g_\infty>2\varepsilon$.}
\end{proposition}
\begin{proof}
\Lang{Every protocol belongs to some finite budget, proving both inequalities in \eqref{eq:sup}. A positive supremum contains a positive witness. For a fixed protocol, the two possible observation intervals are disjoint exactly when $|r_p(W)-r_p(V)|>2\varepsilon$. The definition of supremum supplies a witness above any strictly smaller threshold. At equality the intervals can touch, and if $g_\infty=2\varepsilon$ no protocol exceeds the threshold.}{Каждый протокол входит в некоторый конечный бюджет, что доказывает обе стороны \eqref{eq:sup}. Положительный супремум содержит положительное свидетельство. Для фиксированного протокола два возможных интервала наблюдений не пересекаются ровно при $|r_p(W)-r_p(V)|>2\varepsilon$. Определение супремума даёт свидетельство выше любого строго меньшего порога. При равенстве интервалы могут касаться; если $g_\infty=2\varepsilon$, ни один протокол не превышает порог.}
\end{proof}
\Lang{This is a worst-case bounded-error statement, not a statistical power theorem. Independent repeated noise, correlations, confidence levels and adaptive stopping need their own model. Tolerance indistinguishability is not generally transitive, as already noted in WR-I: responses $0,\varepsilon,2\varepsilon$ with tolerance $\varepsilon>0$ give a counterexample. The compact finite-separation argument belongs to WR-I and is not repeated here.}{Это утверждение об ограниченной ошибке в худшем случае, а не теорема о статистической мощности. Независимые повторения шума, корреляции, уровни доверия и адаптивная остановка требуют своей модели. Неразличимость в пределах допуска в общем случае нетранзитивна, как уже отмечено в WR-I: отклики $0,\varepsilon,2\varepsilon$ при допуске $\varepsilon>0$ дают контрпример. Аргумент конечного разделения на компактной области принадлежит WR-I и здесь не повторяется.}

\begin{proposition}[\Lang{Finite adaptive closure}{Конечное адаптивное замыкание}]\label{prop:adaptive}
\Lang{Fix a finite-depth protocol tree with finite outcome alphabets. Actions and stopping decisions are chosen by one policy whose randomization is independent of the candidate world. Feasible actions and resource certificates are common at each history. If for $W,V$ every common feasible action $p$ at every possible history $h$ has the same conditional outcome law $q_W(y\mid h,p)=q_V(y\mid h,p)$, then the full stopped transcript laws agree.}{Зафиксируем дерево протокола конечной глубины с конечными алфавитами исходов. Действия и остановку выбирает одна политика, рандомизация которой независима от мира-кандидата. Допустимые действия и ресурсные сертификаты общие при каждой истории. Если для $W,V$ каждое общее допустимое действие $p$ при каждой возможной истории $h$ имеет одинаковый условный закон исходов $q_W(y\mid h,p)=q_V(y\mid h,p)$, то законы полной записи с остановкой совпадают.}
\end{proposition}
\begin{proof}
\Lang{The empty-history probability agrees. Inductively a child-history probability is the common parent probability times the common policy action probability times the common conditional response probability. Stopping adds a common factor and can be represented by an absorbing label. Induction through the finite depth proves equality of every transcript probability.}{Вероятность пустой истории одинакова. По индукции вероятность дочерней истории равна общей вероятности родителя, умноженной на общую вероятность выбора действия политикой и на общую условную вероятность отклика. Остановка добавляет общий множитель и представляется поглощающей меткой. Индукция по конечной глубине доказывает равенство вероятностей каждой записи.}
\end{proof}
\Lang{Equality of isolated one-step marginals is insufficient: two bits can each be fair in both worlds but always equal in one and always opposite in the other. A two-step transcript distinguishes them. Likewise, two individually affordable actions need not be jointly affordable. The proposition requires conditional laws and costs of complete policies; it does not justify enlarging a static family without an audit.}{Равенства отдельных одношаговых распределений недостаточно: два бита могут быть равновероятными в обоих мирах, но всегда совпадать в одном и всегда различаться в другом. Двухшаговая запись их различает. Кроме того, два доступных по отдельности действия могут быть недоступными вместе. Утверждение требует условных законов и стоимости полных политик; оно не оправдывает расширение статического семейства без проверки.}

\begin{example}[\Lang{Pairwise access is not uniform access}{Попарный доступ не означает равномерного доступа}]\label{ex:noncompact}
\Lang{On $C=[0,\infty)$ use $r_T(\theta)=(T-\theta)_+$ for finite $T\ge0$, in dimensionless toy units. Every two distinct parameters have a separating finite $T$. For every fixed finite deadline $T_*$, however, all $\theta\ge T_*$ give zero responses for all $T\le T_*$. Thus no finite deadline identifies the whole noncompact carrier. This does not contradict WR-I's compact-domain result.}{На $C=[0,\infty)$ положим $r_T(\theta)=(T-\theta)_+$ для конечного $T\ge0$ в безразмерных единицах учебной модели. Любые два различных параметра различаются при некотором конечном $T$. Однако при каждом фиксированном конечном сроке $T_*$ все $\theta\ge T_*$ дают нулевые отклики для всех $T\le T_*$. Поэтому никакой конечный срок не идентифицирует весь некомпактный носитель. Это не противоречит результату WR-I на компактной области.}
\end{example}

\section{\Lang{A causal winding probe}{Причинный протокол обхода тора}}\label{sec:winding}
\Lang{Use a prescribed flat FLRW metric with $a\in C^1([0,\infty))$, $a>0$, $a(0)=1$:}{Используем заданную плоскую метрику FLRW с $a\in C^1([0,\infty))$, $a>0$, $a(0)=1$:}
\begin{equation}\label{eq:metric}
 ds^2=-c^2dt^2+a(t)^2d\boldsymbol x^2,\qquad
 S(T)=c\int_0^T\frac{dt}{a(t)},\qquad S_\infty=\lim_{T\to\infty}S(T).
\end{equation}
\Lang{Coordinates $\boldsymbol x$ and $S$ have length units; $a$ is dimensionless and $c$ has units length/time. Compare spatial $\mathbb R^3$ with $\mathbb T_L^3=\mathbb R^3/(L\mathbb Z)^3$, $L>0$, under the same $a$. A comoving laboratory at the origin launches an ideal tagged null ray at $t=0$ along a known lattice axis. Only its own first return is read. There are no mirrors, scattering, absorption or false returns. In $\mathbb R^3$ this ray never returns. Axis knowledge, tag recognition and collection are supplied calibration assumptions, not inferred topology.}{Координаты $\boldsymbol x$ и $S$ имеют размерность длины; $a$ безразмерен, $c$ имеет размерность длина/время. Сравним пространства $\mathbb R^3$ и $\mathbb T_L^3=\mathbb R^3/(L\mathbb Z)^3$, $L>0$, при одном $a$. Сопутствующая лаборатория в начале координат запускает идеальный меченый нулевой луч при $t=0$ вдоль известной оси решётки. Считывается только его первое возвращение. Зеркал, рассеяния, поглощения и ложных возвращений нет. В $\mathbb R^3$ этот луч не возвращается. Знание оси, распознавание метки и сбор сигнала \textemdash{} заданные условия калибровки, а не выведенная топология.}

\Lang{The remaining conformal span is the standard comoving event-horizon distance when finite, for this observer and future history; it is not the Hubble radius or the past particle horizon \cite{DavisLineweaver}. Topological multiple images and covering-space constructions are established cosmological tools \cite{Luminet}. We use an explicit ideal active protocol rather than an actual catalogue or CMB inference.}{Оставшийся конформный путь при его конечности является стандартным сопутствующим расстоянием до горизонта событий для этого наблюдателя и будущей истории; он не равен радиусу Хаббла или прошлому горизонту частиц \cite{DavisLineweaver}. Топологические кратные изображения и конструкции накрытия \textemdash{} известные космологические инструменты \cite{Luminet}. Мы используем явный идеальный активный протокол, а не анализ реального каталога или реликтового излучения.}

\begin{theorem}[\Lang{Strict first-return criterion}{Строгий критерий первого возвращения}]\label{thm:return}
\Lang{The torus ray has a unique finite first-return time $T_L$ if and only if $L<S_\infty$. When it exists, $S(T_L)=L$. If $L\ge S_\infty$, it gives the same no-return response as the Euclidean ray at every finite deadline.}{Луч на торе имеет единственное конечное время первого возвращения $T_L$ тогда и только тогда, когда $L<S_\infty$. При существовании $S(T_L)=L$. Если $L\ge S_\infty$, при каждом конечном сроке его отклик отсутствия возвращения совпадает с евклидовым.}
\end{theorem}
\begin{proof}
\Lang{The null equation gives $dx/dt=c/a(t)$ on the lifted axis, hence $x(T)=S(T)$. The first positive lattice point is $L$. Since $S'(T)=c/a(T)>0$, $S$ is continuous and strictly increasing from zero. It reaches every $0<L<S_\infty$ exactly once. For every finite $T$, the positive tail integral implies $S(T)<S_\infty$ when the limit is finite; if it is infinite the conclusion is immediate. Equality $L=S_\infty$ therefore gives no finite return.}{Нулевое условие даёт $dx/dt=c/a(t)$ на поднятой оси, откуда $x(T)=S(T)$. Первая положительная точка решётки равна $L$. Поскольку $S'(T)=c/a(T)>0$, функция $S$ непрерывна и строго возрастает от нуля. Она достигает каждого $0<L<S_\infty$ ровно один раз. Для каждого конечного $T$ положительный хвост интеграла даёт $S(T)<S_\infty$, если предел конечен; при бесконечном пределе вывод очевиден. Поэтому равенство $L=S_\infty$ не даёт конечного возвращения.}
\end{proof}
\Lang{The length is $L$, not the injectivity radius $L/2$ and not the $2R$ patch condition of WR-VIII. Arbitrary directions need not close; an unknown orientation requires additional protocols. No-return equivalence here is relative to the declared tagged winding observations. It does not prove agreement of all global classical fields, mode spectra, quantum states or imaginable experiments. In particular, no-return at one deadline does not distinguish a large torus from its cover.}{Длина равна $L$, а не радиусу инъективности $L/2$ и не условию участка $2R$ из WR-VIII. Произвольные направления могут не замыкаться; неизвестная ориентация требует дополнительных протоколов. Эквивалентность отсутствия возвращения относится к объявленным наблюдениям меченого обхода. Она не доказывает совпадения всех глобальных классических полей, спектров мод, квантовых состояний или мыслимых экспериментов. В частности, отсутствие возвращения к одному сроку не различает большой тор и его накрытие.}

\section{\Lang{A time--energy detection frontier}{Граница обнаружения по времени и энергии}}\label{sec:energy}
\Lang{Now impose $a$ nondecreasing, a launch cap $E_{\max}\ge0$, and a detector threshold $E_d>0$, both in joules. The declared geometric-optics packet model conserves photon number and gives collected energy $E_{\rm arr}=E_0/a(T_L)$ for a launch energy $E_0$ at $a(0)=1$, using the standard cosmological redshift relation \cite{DavisLineweaver}. Collection of the entire packet, no losses and a deterministic inclusive detector threshold are idealizations. The cap excludes preparation and laboratory overhead; it is not a universal total-energy bound. No diffraction, backreaction, quantum detection model or realizable apparatus is certified.}{Теперь потребуем неубывания $a$ и введём предел энергии запуска $E_{\max}\ge0$ и порог детектора $E_d>0$ в джоулях. Объявленная модель пакета в геометрической оптике сохраняет число фотонов и даёт собранную энергию $E_{\rm arr}=E_0/a(T_L)$ при энергии запуска $E_0$ и $a(0)=1$ по стандартному соотношению космологического красного смещения \cite{DavisLineweaver}. Сбор всего пакета, отсутствие потерь и детерминированный включающий порог \textemdash{} идеализации. Предел не учитывает подготовку и расходы лаборатории и не является универсальной границей полной энергии. Дифракция, обратное влияние на метрику, квантовая модель обнаружения и реализуемый прибор не удостоверяются.}
\begin{proposition}[\Lang{Conditional resource certificate}{Условный ресурсный сертификат}]\label{prop:resource}
\Lang{There exists a detectable first return by deadline $\tau<\infty$ with $0\le E_0\le E_{\max}$ if and only if}{Обнаружимое первое возвращение к сроку $\tau<\infty$ с $0\le E_0\le E_{\max}$ существует тогда и только тогда, когда}
\begin{equation}\label{eq:certificate}
 L<S_\infty,\qquad T_L\le\tau,\qquad E_{\max}\ge E_{\rm req}(L):=E_d a(T_L).
\end{equation}
\end{proposition}
\begin{proof}
\Lang{Theorem \ref{thm:return} supplies the unique return. Detection requires}{Теорема \ref{thm:return} даёт единственное возвращение. Условие обнаружения имеет вид}
\[
 E_0/a(T_L)\ge E_d\quad\Longleftrightarrow\quad E_0\ge E_d a(T_L).
\]
\Lang{Such an energy lies under the cap exactly under \eqref{eq:certificate}; choose $E_0=E_{\rm req}$. Both finite-deadline and detector equalities are included by their stated conventions.}{Такая энергия помещается под предел ровно при \eqref{eq:certificate}; выбираем $E_0=E_{\rm req}$. Равенства для конечного срока и порога детектора включены согласно заданным соглашениям.}
\end{proof}

\begin{corollary}[\Lang{Exponential benchmark}{Экспоненциальный пример}]\label{cor:exponential}
\Lang{For $a(t)=e^{Ht}$, $H>0$ in inverse seconds, set $x=HL/c$. For $0<x<1$,}{Для $a(t)=e^{Ht}$, $H>0$ в обратных секундах, положим $x=HL/c$. При $0<x<1$}
\begin{equation}\label{eq:exponential}
 S(T)=\frac cH(1-e^{-HT}),\quad S_\infty=\frac cH,\quad
 T_L=-\frac1H\log(1-x),\quad E_{\rm req}=\frac{E_d}{1-x}.
\end{equation}
\Lang{For $E_{\max}\ge E_d$, the detectable positive periods at finite budget are exactly}{При $E_{\max}\ge E_d$ обнаружимые положительные периоды при конечном бюджете имеют вид}
\begin{equation}\label{eq:frontier}
 0<L\le L_{\rm det}(\tau,E_{\max})=
 \frac cH\min\{1-e^{-H\tau},\ 1-E_d/E_{\max}\}.
\end{equation}
\Lang{The set is empty for $E_{\max}<E_d$, and also when the right endpoint is zero.}{При $E_{\max}<E_d$ множество пусто; оно пусто и при нулевой правой границе.}
\end{corollary}
\begin{proof}
\Lang{Integrate \eqref{eq:metric} and solve $S(T_L)=L$. Then $a(T_L)=(1-x)^{-1}$. Substituting into \eqref{eq:certificate} gives $x\le1-e^{-H\tau}$ and $x\le1-E_d/E_{\max}$. Nondecreasing $a\ge1$ excludes detection when $E_{\max}<E_d$.}{Интегрируем \eqref{eq:metric} и решаем $S(T_L)=L$. Затем $a(T_L)=(1-x)^{-1}$. Подстановка в \eqref{eq:certificate} даёт $x\le1-e^{-H\tau}$ и $x\le1-E_d/E_{\max}$. Неубывание $a\ge1$ исключает обнаружение при $E_{\max}<E_d$.}
\end{proof}
\Lang{For $H\tau=\log2$ and $E_{\max}=2E_d$, the endpoint $L=c/(2H)$ is detected exactly at both equalities. At $L=c/H$ the return never occurs at finite time, even with unlimited launch energy. As $x\uparrow1$, $T_L$ diverges logarithmically and $E_{\rm req}$ diverges as $(1-x)^{-1}$. A fixed energy cap creates a stricter detection limit than causal reach. These are different reasons for identical no-detection responses.}{При $H\tau=\log2$ и $E_{\max}=2E_d$ крайний период $L=c/(2H)$ обнаруживается ровно при обоих равенствах. При $L=c/H$ возвращения в конечное время нет даже при неограниченной энергии запуска. При $x\uparrow1$ время $T_L$ расходится логарифмически, а $E_{\rm req}$ как $(1-x)^{-1}$. Фиксированный энергетический предел даёт более строгую границу обнаружения, чем причинная доступность. Это разные причины одинаковых откликов отсутствия обнаружения.}
\begin{figure}[ht]
\centering\includegraphics[width=0.92\linewidth]{figures/frontier\wrehFigureSuffix.pdf}
\caption{\Lang{Synthetic exponential model: normalized return time and required energy. The causal endpoint $x=1$ is excluded; finite resource thresholds may be attained. No cosmological data are fitted.}{Синтетическая экспоненциальная модель: нормированные время возвращения и необходимая энергия. Причинная граница $x=1$ исключена; конечные ресурсные пороги могут достигаться. Подгонки космологических данных нет.}}\label{fig:frontier}
\end{figure}

\section{\Lang{Current compatibility does not fix future reach}{Нынешняя совместимость не фиксирует будущую доступность}}\label{sec:futures}
\begin{theorem}[\Lang{Smooth agreement with different tail convergence}{Гладкое совпадение при различной сходимости хвостов}]\label{thm:tails}
\Lang{Fix $H>0$, $T_*>0$ and $\delta>0$. There are positive, strictly expanding $C^\infty$ scale factors on $\mathbb R$ which both equal $e^{Ht}$ for every $t\le T_*$, but have finite versus infinite $S_\infty$ from $t=0$.}{Зафиксируем $H>0$, $T_*>0$ и $\delta>0$. Существуют положительные, строго расширяющиеся $C^\infty$ масштабные факторы на $\mathbb R$, оба равные $e^{Ht}$ при всех $t\le T_*$, но имеющие конечный либо бесконечный $S_\infty$ от $t=0$.}
\end{theorem}
\begin{proof}
\Lang{Let $\rho(u)=e^{-1/u}$ for $u>0$ and zero otherwise, and set}{Положим $\rho(u)=e^{-1/u}$ при $u>0$ и ноль иначе; введём}
\begin{equation}\label{eq:switch}
 s(u)=\frac{\rho(u)}{\rho(u)+\rho(1-u)},\qquad
 h_-(t)=\begin{cases}H,&t\le T_*,\\
 H[1-s((t-T_*)/\delta)]+\dfrac{Hs((t-T_*)/\delta)}{1+H(t-T_*)},&t>T_*.
 \end{cases}
\end{equation}
\Lang{The denominator defining $s$ is positive everywhere; $s=0$ on $u\le0$, $s=1$ on $u\ge1$, and it is smooth and flat at both join points: every derivative of $e^{-1/u}$ is a polynomial in $1/u$ times $e^{-1/u}$ and tends to zero as $u\downarrow0$. Hence $h_-$ is smooth and positive, agrees with $H$ through $T_*$, and equals $H/[1+H(t-T_*)]$ after $T_*+\delta$. Define}{Знаменатель в определении $s$ всюду положителен; $s=0$ при $u\le0$, $s=1$ при $u\ge1$, функция гладка и плоска в обеих точках склейки: каждая производная $e^{-1/u}$ является многочленом от $1/u$, умноженным на $e^{-1/u}$, и стремится к нулю при $u\downarrow0$. Поэтому $h_-$ гладок и положителен, совпадает с $H$ до $T_*$ включительно и равен $H/[1+H(t-T_*)]$ после $T_*+\delta$. Определим}
\begin{equation}\label{eq:tails}
 a_+(t)=e^{Ht},\qquad a_-(t)=\exp\!\left(\int_0^t h_-(v)\,dv\right).
\end{equation}
\Lang{Both have the stated past and $\dot a_\pm/a_\pm>0$. For $t\ge T_*+\delta$,}{Оба имеют указанное прошлое и $\dot a_\pm/a_\pm>0$. При $t\ge T_*+\delta$}
\begin{equation}\label{eq:linear-tail}
 a_-(t)=A[1+H(t-T_*)],\qquad
 A=\frac{a_-(T_*+\delta)}{1+H\delta}>0.
\end{equation}
\Lang{Its reciprocal has a divergent logarithmic tail integral, whereas $a_+$ gives $S_\infty=c/H$. Smoothness also gives agreement of all derivatives at $T_*$.}{Интеграл обратной величины имеет расходящийся логарифмический хвост, тогда как $a_+$ даёт $S_\infty=c/H$. Гладкость также обеспечивает совпадение всех производных при $T_*$.}
\end{proof}
\Lang{In the calibrated flat current carrier, both histories have $H(z)=H$ and $D(z)=cz/H$ on any supplied finite $[0,Z]$. A period $L\ge c/H$ has no finite winding return under $a_+$ but does return under $a_-$, at some finite time and finite threshold energy. To compare topology lifts on the restricted current patch of WR-VIII, also impose $L>2cZ/H$; choosing $L>\max\{c/H,2cZ/H\}$ meets both requirements. Equality of other current fields still requires the explicit patch-copy premise. Until $T_*$ the backgrounds agree exactly. This is a change of future tail, not a compact bump. A prescribed matter evolution law, analyticity, or a fixed cosmological parameter family can exclude this construction; none is supplied here. It is a kinematic counterexample to inferring permanent access from current background responses, not a pair of solutions of one declared matter theory.}{В калиброванном плоском нынешнем носителе обе истории имеют $H(z)=H$ и $D(z)=cz/H$ на любом заданном конечном $[0,Z]$. Период $L\ge c/H$ не даёт конечного возвращения при $a_+$, но даёт его при $a_-$ в некоторое конечное время и при конечной пороговой энергии. Для сравнения топологических продолжений на ограниченном нынешнем участке WR-VIII дополнительно требуется $L>2cZ/H$; выбор $L>\max\{c/H,2cZ/H\}$ удовлетворяет обоим условиям. Совпадение остальных нынешних полей по-прежнему требует явной предпосылки копирования участка. До $T_*$ фоны точно совпадают. Здесь меняется хвост будущего, а не компактный всплеск. Предписанный закон эволюции материи, аналитичность или фиксированное семейство космологических параметров могут исключить эту конструкцию; здесь они не заданы. Это кинематический контрпример к выводу о постоянном доступе из нынешних фоновых откликов, а не пара решений одной объявленной теории материи.}

\begin{proposition}[\Lang{Compact future changes preserve convergence type}{Компактные изменения будущего сохраняют тип сходимости}]\label{prop:compact-tail}
\Lang{If positive continuous $a,b$ agree for every $t\ge T_1<\infty$, then $\int_0^\infty dt/a(t)$ is finite if and only if $\int_0^\infty dt/b(t)$ is finite.}{Если положительные непрерывные $a,b$ совпадают при всех $t\ge T_1<\infty$, то $\int_0^\infty dt/a(t)$ конечен тогда и только тогда, когда конечен $\int_0^\infty dt/b(t)$.}
\end{proposition}
\begin{proof}
\Lang{Each reciprocal is bounded on the compact interval $[0,T_1]$ because its denominator has a positive minimum. The remaining tails are identical.}{Каждая обратная величина ограничена на компактном интервале $[0,T_1]$, поскольку знаменатель имеет положительный минимум. Остальные хвосты совпадают.}
\end{proof}
\Lang{The numerical value of a finite span can still change; a fixed $L$ can cross that value. The compact future bumps of WR-VIII therefore establish temporal ambiguity but cannot alone switch finite to infinite conformal reach. Theorem \ref{thm:tails} supplies the missing noncompact-tail construction.}{Численное значение конечного пути всё же может меняться; фиксированный $L$ может пересечь это значение. Поэтому компактные будущие всплески WR-VIII устанавливают временную неоднозначность, но сами по себе не могут переключить конечную конформную доступность в бесконечную. Теорема \ref{thm:tails} даёт необходимую конструкцию некомпактного хвоста.}
\begin{figure}[ht]
\centering\includegraphics[width=0.92\linewidth]{figures/tails\wrehFigureSuffix.pdf}
\caption{\Lang{Synthetic $H=c=1$, $T_*=1$, $\delta=1/2$. The histories and reach agree until the switch; thereafter the exponential span approaches one while the smooth linear-tail span is unbounded. Finite plot values do not prove divergence; Eq. \eqref{eq:linear-tail} does.}{Синтетический пример $H=c=1$, $T_*=1$, $\delta=1/2$. Истории и доступность совпадают до переключения; затем экспоненциальный путь стремится к единице, а путь с гладким линейным хвостом неограничен. Конечные значения графика не доказывают расходимость; её доказывает формула \eqref{eq:linear-tail}.}}\label{fig:tails}
\end{figure}

\FloatBarrier
\section{\Lang{Unknown futures and the order of quantifiers}{Неизвестные будущие и порядок кванторов}}\label{sec:quantifiers}
\begin{proposition}[\Lang{Possible versus excluded winding access}{Возможный и исключённый доступ через обход}]\label{prop:robust}
\Lang{Let $\mathcal G\ne\varnothing$ be a declared family of positive future scale factors as in Section \ref{sec:winding}. A fixed period $L>0$ has no finite causal winding return for any $\gamma\in\mathcal G$ if and only if}{Пусть $\mathcal G\ne\varnothing$ \textemdash{} объявленное семейство положительных будущих масштабных факторов из раздела \ref{sec:winding}. Фиксированный период $L>0$ не имеет конечного причинного возвращения ни при одном $\gamma\in\mathcal G$ тогда и только тогда, когда}
\begin{equation}\label{eq:robust}
 L\ge \sup_{\gamma\in\mathcal G}S_\infty^{(\gamma)}.
\end{equation}
\Lang{A return is possible in at least one future exactly when some $S_\infty^{(\gamma)}>L$. Return in every future requires $L<S_\infty^{(\gamma)}$ for every $\gamma$; it does not imply a common finite resource certificate.}{Возвращение возможно хотя бы в одном будущем ровно тогда, когда некоторое $S_\infty^{(\gamma)}>L$. Возвращение в каждом будущем требует $L<S_\infty^{(\gamma)}$ для каждого $\gamma$; это не влечёт общего конечного ресурсного сертификата.}
\end{proposition}
\begin{proof}
\Lang{Apply Theorem \ref{thm:return} to each future. Having $L\ge S_\infty^{(\gamma)}$ for all $\gamma$ is exactly \eqref{eq:robust}, with the usual extended supremum. The last claim follows from the example below.}{Применим теорему \ref{thm:return} к каждому будущему. Условие $L\ge S_\infty^{(\gamma)}$ для всех $\gamma$ равносильно \eqref{eq:robust} с обычным расширенным супремумом. Последний вывод следует из примера ниже.}
\end{proof}
\begin{example}[\Lang{All returns exist, but no uniform budget}{Все возвращения существуют, но общего бюджета нет}]\label{ex:nonuniform}
\Lang{Take $H_n=H_0(1-1/n)$, $a_n=e^{H_nt}$, $n\ge2$, and $L=c/H_0$. Every history has $S_\infty^{(n)}=L/(1-1/n)>L$, while their infimum equals $L$ and is not attained. Equations \eqref{eq:exponential} give}{Положим $H_n=H_0(1-1/n)$, $a_n=e^{H_nt}$, $n\ge2$ и $L=c/H_0$. Каждая история имеет $S_\infty^{(n)}=L/(1-1/n)>L$, но их инфимум равен $L$ и не достигается. Формулы \eqref{eq:exponential} дают}
\begin{equation}\label{eq:nonuniform}
 T_n=\frac{\log n}{H_0(1-1/n)}\longrightarrow\infty,\qquad
 E_{{\rm req},n}=nE_d\longrightarrow\infty.
\end{equation}
\Lang{Thus $\forall\gamma\,\exists(\tau,E)$ cannot be replaced by $\exists(\tau,E)\,\forall\gamma$. Requiring $L<\inf_\gamma S_\infty^{(\gamma)}$ would incorrectly exclude this example. These $H_n$ are a family with uncertain current rate, not the exactly shared current data of Theorem \ref{thm:tails}.}{Поэтому $\forall\gamma\,\exists(\tau,E)$ нельзя заменить на $\exists(\tau,E)\,\forall\gamma$. Требование $L<\inf_\gamma S_\infty^{(\gamma)}$ ошибочно исключило бы этот пример. Здесь $H_n$ образуют семейство с неопределённой нынешней скоростью, а не точно общие нынешние данные теоремы \ref{thm:tails}.}
\end{example}

\section{\Lang{Scientific ceiling and closure of the cycle}{Научная граница утверждений и завершение цикла}}\label{sec:ceiling}
\Lang{The central distinction is between a response class at current resources and a horizon relative to a justified future access record. In the benchmark, the record explicitly fixes future expansion, lattice orientation, packet transport, detector threshold and allowed launches. Changing any of these can change the certificate. A causal horizon restricts information propagation; it does not say that no mathematical or physical structure exists beyond it. Exact equality of the winding response says nothing about unrelated protocols.}{Главное различие проходит между классом откликов при нынешних ресурсах и горизонтом относительно обоснованных будущих условий доступа. В примере явно фиксируются будущее расширение, ориентация решётки, перенос пакета, порог детектора и допустимые запуски. Изменение любого из этих условий может менять сертификат. Причинный горизонт ограничивает перенос информации; он не означает отсутствия математической или физической структуры за ним. Точное равенство отклика обхода ничего не говорит о других протоколах.}

\Lang{A mathematical FLRW history is not automatically a solution of a prescribed matter theory or a realistic universe. A usable physical analysis must add a well-posed matter/gravity model, perturbations and fields, initial and boundary data, observer survival, apparatus dynamics, directional uncertainty, losses and statistical detection. The current article certifies none of these. Known event-horizon and redshift formulas, torus coverings, bounded-error intervals and conditional-product arguments are prior tools. The contribution is their explicit WREH access/resource contract and the future-tail counterexample; broader priority remains unestablished.}{Математическая история FLRW не становится автоматически решением предписанной теории материи или реалистичной вселенной. Для физического анализа нужны корректно поставленная модель материи и гравитации, возмущения и поля, начальные и граничные данные, существование наблюдателя, динамика приборов, неопределённость направления, потери и статистика обнаружения. Здесь это не удостоверяется. Известные формулы горизонта событий и красного смещения, накрытия тора, интервалы ограниченной ошибки и аргументы условных произведений \textemdash{} предшествующие инструменты. Вклад состоит в явных условиях доступа и ресурсов WREH и контрпримере с хвостом будущего; более широкий приоритет не установлен.}

\Lang{This is the ninth and final technical article of the manifesto's first cycle. It does not declare all nine articles published, reviewed or merged. WR-V remains a separate review draft; WR-VII publication is author-reported pending deposited-file verification; WR-VIII v0.2 is deposited at DOI \nolinkurl{10.5281/zenodo.23266980}. Earlier frozen sources are retained. Open PR dependencies and scientific review remain separate from the completion of this draft. A later synthesis requires its own scope and evidence. WREH and Boundary Compensation remain distinct programmes.}{Это девятая и последняя техническая статья первого цикла манифеста. Она не объявляет все девять статей опубликованными, отрецензированными или объединёнными в репозитории. WR-V остаётся отдельным черновиком для рецензирования; публикация WR-VII сообщена автором и ожидает проверки депонированных файлов; WR-VIII v0.2 депонирована с DOI \nolinkurl{10.5281/zenodo.23266980}. Ранние замороженные источники сохранены. Открытые зависимости PR и научное рецензирование отделены от завершения этого черновика. Последующий синтез требует собственной области и доказательств. WREH и Boundary Compensation остаются отдельными программами.}

\paragraph{\Lang{Reproduction and audit}{Воспроизведение и аудит}.}
\Lang{The package contains shared bilingual LaTeX, analytic figures, an offline demonstration, symbolic/numerical and adversarial boundary checks, a source map and an audit. Computation checks finite witnesses and implementations; the proofs establish the infinite-tail claims. AI-assisted drafting and automated technical audit are disclosed. They are not independent external scientific peer review, and no journal acceptance or WR-IX deposit is claimed.}{Комплект содержит общий двуязычный LaTeX, аналитические иллюстрации, автономную демонстрацию, символьные и численные проверки, проверку граничных случаев, карту источников и аудит. Вычисления проверяют конечные свидетельства и реализации; утверждения о бесконечных хвостах устанавливаются доказательствами. Раскрываются подготовка текста с помощью ИИ и автоматизированный технический аудит. Они не являются независимым внешним научным рецензированием; принятие журналом или депонирование WR-IX не заявляются.}

\begin{thebibliography}{99}
\raggedright
\bibitem{Manifesto} A. A. Malachevsky. \emph{WREH-00: Manifesto and Research Programme}, v0.1 (2026), \Lang{Sections 8, 10--11}{Разделы 8, 10--11}. \url{https://github.com/AIDevelopersMonster/WREH/tree/3963c3b71f464b7de85750fd4a51d680943d5162/docs/WREH-00}.
\bibitem{WRI} A. A. Malachevsky. \emph{Response-Defined World Spaces}, WR-I v0.3 (2026). DOI: \texttt{10.5281/zenodo.23210258}.
\bibitem{WRII} A. A. Malachevsky. \emph{Finite Consistency and Global World Realizability}, WR-II v0.5 (2026). DOI: \texttt{10.5281/zenodo.23224154}.
\bibitem{WRIII} A. A. Malachevsky. \emph{Geometry of Admissible World Fibres}, WR-III v0.6 (2026). DOI: \texttt{10.5281/zenodo.23228682}.
\bibitem{WRIV} A. A. Malachevsky. \emph{Response-Conditioned Global Completion and the Status of the Past}, WR-IV v0.6 (2026). DOI: \texttt{10.5281/zenodo.23248138}.
\bibitem{WRV} A. A. Malachevsky. \emph{Experiment and Realization Selection}, WR-V v0.3 (2026). \Lang{Unpublished repository draft.}{Неопубликованный черновик репозитория.}
\bibitem{WRVI} A. A. Malachevsky. \emph{The Aquarium Bounds}, WR-VI v0.2 (2026). DOI: \texttt{10.5281/zenodo.23254619}.
\bibitem{WRVII} A. A. Malachevsky. \emph{Energy-Time Frontiers of World Realizability}, WR-VII v0.2 (2026). \Lang{Author-reported DOI}{Сообщённый автором DOI}: \texttt{10.5281/zenodo.23264253}; \Lang{deposited files pending verification}{депонированные файлы ожидают проверки}.
\bibitem{WRVIII} A. A. Malachevsky. \emph{The Cosmological Fibre: Universes Compatible with Our Universe Today}, WR-VIII v0.2 (2026). DOI: \nolinkurl{10.5281/zenodo.23266980}.
\bibitem{DavisLineweaver} T. M. Davis and C. H. Lineweaver. \emph{Expanding Confusion: Common Misconceptions of Cosmological Horizons and the Superluminal Expansion of the Universe}. Publ. Astron. Soc. Aust. \textbf{21}, 97--109 (2004). DOI: \texttt{10.1071/AS03040}. \texttt{arXiv:astro-ph/0310808v2}, \Lang{Appendix A, Eqs. (23), (28)}{Приложение A, формулы (23), (28)}.
\bibitem{Luminet} J.-P. Luminet. \emph{The Status of Cosmic Topology after Planck Data}. Universe \textbf{2}(1), 1 (2016). DOI: \texttt{10.3390/universe2010001}. \texttt{arXiv:1601.03884v2}, \Lang{Sections 2--3}{Разделы 2--3}.
\end{thebibliography}
